\documentclass[UTF8, a4paper, zihao=-4]{ctexart}
% 模板：ctexart + natbib 数字引用（EXECUTION_PLAN §8.3；课程模板下发后迁移）
\usepackage[numbers, sort&compress]{natbib}
\usepackage{geometry}
\geometry{margin=2.5cm}
\usepackage{booktabs}
\usepackage{longtable}
\usepackage{graphicx}
\usepackage{url}
\usepackage{tikz}
\usetikzlibrary{arrows.meta, positioning}
\usepackage{hyperref}

\title{\textbf{大模型智能体的长期记忆：\\组织、更新与可靠性评测综述}}
\author{成员 A\quad 成员 B\quad 成员 C}  % 待真人姓名替换
\date{2026 年 10 月}

\begin{document}
\maketitle

\begin{abstract}
大语言模型智能体在跨会话、长期任务中需要外部记忆以保留经验、跟踪状态并复用知识。本文综述 LLM 智能体长期记忆的组织、更新与可靠性评测：基于三源检索与两阶段筛选（691 项候选经题名摘要初筛与正文深读，40 篇核心阅读池、39 项进入正文引用），从生命周期视角梳理研究脉络。主要发现：（1）组织与检索机制的增益高度依赖任务与设置——图结构记忆在 LongMemEval 上的 token F1（0.417）低于匹配预算的扁平向量基线（0.468），结构化组织并非必然优于扁平存储；（2）遗忘与更新是低成本可行操作（一次图剪枝移除 9.8\% 节点而 token F1 变化不显著），但存在信息损失与误删权衡；（3）错误写入、摘要失真、错误检索与错误使用构成独立于恶意攻击的可靠性问题，且幻觉在多阶段流水线中沿传播链累积；（4）评测从会话长度转向能力维度（知识更新、弃答），但``更新正确性''与``记忆何时真正有帮助''仍缺乏统一的成本—收益度量。综述同时给出一个以逐字证据验证与人工核验节点约束幻觉的综述智能体实现及其小规模评估。局限：领域以近期预印本为主，结论的条件性强，人工核验进行中。
\par\textbf{关键词：}大语言模型；智能体；长期记忆；记忆更新；可靠性；评测
\end{abstract}

% intro：生成于阅读卡+证据；引用检查通过
\section{引言}

大语言模型（Large Language Model, LLM）驱动的智能体在长期交互任务中面临一个核心约束：模型的原生上下文窗口是有限且易失的。现有研究表明，即便上下文窗口被扩展，它仅提供会话内的召回能力，并不能自动实现跨会话的持久保留、共享与检索，因此不构成真正意义上的长期记忆（long-term memory, LTM）\cite{W0039}。基于提示的 LLM 智能体受上下文窗口大小限制，无法回忆已见内容，也难以在少量演示之外持续学习\cite{W0618}。在静态检索增强生成（Retrieval-Augmented Generation, RAG）范式下，检索通常针对单一查询执行，缺乏面向跨会话可写状态的读写接口；而参数化知识既难以在推理时被精确更新，也难以携带个体经验与交互轨迹。作者在 MemGPT 中进一步指出，LLM 需要借助外部存储与类似操作系统的自主记忆管理，才能突破固定上下文窗口的约束\cite{W0002}。这些限制共同构成了长期任务引入外部可写记忆的基本动因。

为此，研究者已提出多种长期记忆实现方案。生成式智能体（Generative Agents）以自然语言记忆流记录智能体全部经验，通过相关性、新近性与重要性综合筛选，并引入反思机制将记忆综合为高层推断\cite{W0001}。MemGPT 将主上下文划分为只读系统指令、可写工作上下文与消息队列，外部上下文类比磁盘存储，记忆编辑与检索完全由 LLM 自主完成\cite{W0002}。RET-LLM 旨在提供通用读写记忆，支持跨任务知识存储与检索，且不修改 LLM 架构\cite{W0680}。ExpeL 通过自主收集经验并提取跨任务洞察实现学习与回忆\cite{W0618}。Mem0 与 A-Mem 则分别从生产级记忆操作与动态组织机制角度，尝试构建可扩展、可演化的记忆系统\cite{W0682,W0681}。然而，作者指出当前系统多停留在基础存储与检索层面，即使引入图数据库，也可能因固定操作与预定 schema 而受限\cite{W0681}。同时，术语与边界的不统一增加了横向比较的困难；CoALA 框架从认知架构角度区分了程序性、语义与情景三类记忆，为理解智能体内部状态提供了概念参照\cite{W0140}。安全综述则将长期记忆界定为可写、跨会话的持久记忆，并提出涵盖写入、存储、检索、执行、共享与传播、遗忘与回滚的生命周期框架\cite{W0039}。

在上述背景下，本文关注长期任务中外部记忆机制的读写闭环及其对智能体行为的影响，而非参数化记忆更新，也不研究单查询检索式 RAG 的检索质量本身。围绕跨会话可写外部记忆，本文提出以下研究问题：

\textbf{RQ1（定义与边界）：}长期记忆包含哪些信息和机制？它与上下文历史、静态知识检索、参数知识有什么关系？

\textbf{RQ2（组织与使用）：}写入、组织和检索方法如何影响信息可用性、任务表现和成本？

\textbf{RQ3（更新与遗忘）：}如何处理信息变化、时间条件、冲突、过时记忆和必要的删除？

\textbf{RQ4（可靠性）：}错误摘要、错误反思或来源混淆如何进入记忆并影响后续决策？现有方法怎样发现和控制这些问题？

\textbf{RQ5（评测）：}基准和实验是否真正考察跨会话记忆、更新能力和长期任务表现？不同结果能否公平比较？

上述研究问题限定在跨会话可写的外部记忆范围之内。本文综合已有工作，但不将参数化记忆纳入讨论，也不对单查询检索式 RAG 的索引与排序算法作系统比较。对于跨系统、跨任务的可比较度量与结论，若素材未直接给出，本文仅作定性综合，并明确标注为【待核验】。

% concepts：内容待 G4/G5 基于证据生成；每条核心陈述将关联 claim_id 注释。
\section{概念界定与综述方法}

（本节内容待生成。）

% taxonomy：内容待 G4/G5 基于证据生成；每条核心陈述将关联 claim_id 注释。
\section{研究脉络与方法分类}

（本节内容待生成。）

% update_reliability：生成于阅读卡+证据；引用检查通过


记忆系统的长期可靠性首先取决于更新与遗忘机制能否在不损害有用信息的前提下淘汰过时或冗余内容。现有工作从选择性遗忘、治理原语、存储/使用分离、再巩固、版本规范化与偏好更新等不同路径探索这一问题。

图结构记忆中的选择性遗忘是代表性做法。W0034 在 LongMemEval 基准上以提取式流水线评测图结构记忆（27021 个节点，候选生成预算为 5 个检索根），其选择性遗忘模块按 recency、访问频率、度中心性和 age 加权得分剪枝低分节点。作者报告，一次应用遗忘移除 9.8% 节点与 9.5% 存储字节，token F1 变化 +0.001（95% CI [-0.015, +0.016]），judged correctness 下降 1.6 点；同时图结构记忆的 token F1 为 0.417，低于扁平向量基线的 0.468（配对 bootstrap Δ=-0.050）\cite{W0034}。这说明在该设定下，图结构与剪枝策略并未带来召回优势，遗忘更多是成本控制手段而非性能提升手段。

与基于启发式得分的剪枝不同，W0199 提出 Memory Worth 作为遗忘治理的在线原语：系统在写入时记录记忆是否属于当前集合，在读取时通过任务成败反馈持续估计条件成功概率。作者称在合成环境（真实效用已知、20 个独立种子、10000 回合）中，Memory Worth 与真实效用的斯皮尔曼秩相关达到 ρ=0.89±0.02，而从不更新的系统为 ρ=0.00；在真实文本微实验（all-MiniLM-L6-v2 嵌入检索、3000 回合）中，过时记忆 MW=0.17 低于低价值阈值，专家记忆 MW=0.77。作者特别说明该度量捕获的是关联而非因果，应视为最小操作原语\cite{W0199}。

W0178 则把遗忘与信任、防投毒统一在同一预算策展框架下。CURATOR 以“净值/字节”单一评分管理经验记忆的保留、共享与信任，净值含价值项与危害项（含来源项）。作者称在 Qwen2.5-3B、Qwen2.5-7B、Phi-3.5-mini 上，仅按价值筛选或全量保留均使注入攻击成功率保持很高，而加入来源项的净值得分可将 ASR 降至零\cite{W0178}。这一结果表明，遗忘决策不能只看信息价值，还需显式建模来源可信度。

关于“该遗忘什么”，W0180 主张将存储与使用分离：归档保留源观察，查询条件视图决定哪些证据被用于回答，遗忘因而是查询局部的——事实可以继续存储但在特定答案中被抑制使用。作者在五个任务套件、四个语言模型骨干上评测 RD-Forget，报告在事实整合任务上相对更强基线的提升为 11 到 26 个百分点不等，且无遗忘配置与查询无条件配置在消融中得分最低\cite{W0180}。

W0012 提出检索驱动的再巩固（REALM）：将插入、合并与遗忘统一为单一整合阶段，把遗忘视为冗余处理。作者称 REALM 在 LoCoMo 上平均准确率 75.97%、在 LongMemEval 上 65.11%，分别比最强基线高 7.17 和 1.31 个百分点\cite{W0012}。这一思路与分离式流水线形成鲜明对比，但证据未详述冲突与信息变化的具体处理。

版本规范化处理的是跨模型记忆漂移问题。W0077 以“解释熵”刻画记忆对象被不同模型赋予不同解读的程度，主张将消歧从检索时前置到写入时。作者在 n=176 合成记忆对象、三个下游模型家族的配对准随机测试中发现，全规范化记忆相对原始记忆在情感一致性上有未校正的显著改善（+0.050，95% bootstrap CI [0.013, 0.086]，配对 t 检验 p=0.0），但多重比较校正后不再显著；事实漂移在任何校正水平下均不显著。作者明确将这些结果定性为探索性而非确认性\cite{W0077}。本文综合认为，版本规范化在概念上切中跨模型部署的痛点，但其经验证据强度尚不足以支撑强结论。

偏好更新方面，W0011 批评 MemGPT、MemoryBank、A-MEM、SCM、MemInsight 等系统默认用户行为静态、缺乏偏好漂移自适应更新，并提出融合短期滑动窗口与长期指数移动平均的偏好表征。作者称 PAMU 在 λ ∈ [0.4, 0.7] 与 β ∈ [0.8, 0.95] 范围内性能稳定\cite{W0011}。这提示偏好记忆需要同时追踪波动与趋势，但该结论的量化评测细节在所给证据中未展开【待核验】。

\subsection{非恶意可靠性}

即使没有恶意攻击，记忆系统在抽取、更新、检索与压缩各环节都可能因信息失真而不可靠。

错误传播是流水线式智能体系统的突出风险。W0688 在 4 智能体 LangGraph 流水线（Researcher→Analyst→Writer→Reviewer）上，以 FinanceBench 数据、346 条注入幻觉、gpt-4o-mini 智能体与 gpt-4o 评判进行评测。作者称幻觉检测率从 Stage 1 的 72.0% 单调下降到 Stage 4 的 50.9%；边界门控使幻觉生存率从 58.4% 降至 16.2%，而末端检查仅改善 2.3 个百分点，作者称该改善在统计上可忽略\cite{W0688}。这一对比说明，在错误产生的边界处设防比在输出末端检查更有效。

记忆幻觉可能在多个操作环节累积。W0261 提出操作级评测 HaluMem，区分记忆抽取、记忆更新与记忆问答三个任务。作者称现有记忆系统倾向于在抽取和更新阶段产生并累积幻觉，随后向问答阶段传播错误；并指出 LoCoMo、LongMemEval、PrefEval、PersonaMem 等基准采用端到端黑盒评估，难以定位幻觉发生的具体环节\cite{W0261}。本文综合认为，操作级诊断对可靠性的归因分析不可或缺，但该结论的经验基础限于 HaluMem-Medium 与 HaluMem-Long 数据集的评测范围。

检索环节的另一类失败是“证据不具可采纳性”。W0690 指出，现有评测只关注召回和最终答案正确性，无法识别检索到正确内容但因语境不匹配而不应被使用的失败模式。作者在 RHELM 与 MemOps 共 3767 个查询的 top-20 设置中报告，采用发布元数据参照后，Top-20 anchor recall 从 0.432 提升到 0.533，80% recall feasibility 从 0.237 提升到 0.311，exact similarity evaluations 下降 98.3%；在冻结 72 案例开发诊断中，仅文本验证器在 1% 必需锚点误拒上限内无法检测违规\cite{W0690}。这表明检索可采纳性需要超出语义相似度的验证信号。

来源追踪是抑制错误传播与幻觉的通用手段。W0010 提出来源感知记忆系统，存储对话事实的同时记录来源与置信度，并在读取时施加引用锁纪律：答案只能引用已打开的证据，无法组成引用集则弃权。作者报告在 LongMemEval 上达 95.60%（较此前最强系统 +0.73 分），在 LoCoMo 上达 93.60%（+1.10 分）；在八个骨干 LLM 的受控实验中，准确率仅差 3.4 分而单查询成本差约 30 倍，作者称能以约 20 倍低成本达到接近 SOTA 质量\cite{W0010}。该工作的证据强度相对较高，覆盖两个基准与多骨干模型。

从信息论视角看，所有紧凑记忆形式——KV 缓存、gist 向量、循环状态与智能体笔记——共享同一 Fano 下界约束。W0207 指出，查询无关压缩的有效每查询预算被查询分布熵 H(Q) 削减，快照式方法在查询分布变化时性能塌陷；作者还指出现有评测仅仔细测量单轮长上下文压缩，智能体实际执行的重复压缩几乎未被测量\cite{W0207}。这为理解记忆压缩失真提供了一个统一的理论参照，但其 COMPACT-Bench 的实证数据在给定证据中未提供【待核验】。

\subsection{恶意攻击与防御及其代价}

当记忆系统暴露于对抗性输入时，选择性提取与改写机制本身可能被利用。W0319 提出的 MemPoison 针对具有选择性提取与改写两阶段的 LLM 智能体记忆流水线，攻击者仅通过交互注入、无直接记忆写权限，但可借助本地白盒代理嵌入模型优化触发器。作者称 MemPoison 跨多智能体领域与记忆机制的攻击成功率最高可达 0.95，优于静态 RAG 或直接写库类基线\cite{W0319}。这揭示了一个深层矛盾：为可靠性设计的筛选流水线反而为绕过滤创造了新攻击面。

防御层面的综合策略体现在 W0076：检索时采用信任感知过滤与排序（含时间衰减），写入时通过复合信任分数门控。该工作以 guard agent 加 memory sanitization 构成防御，模型为 GPT-4o-mini-class 与 Gemini-2.0-Flash，场景为 eICU 与 MIMIC 临床 prompt 集；但给定证据未提供定量结果\cite{W0076}。本文综合判断该方向合理，但防御效果尚缺乏可核验的数字支撑【待核验】。

防御的良性用例代价常被忽视。W0312 在完全良性的 LoCoMo 对话流量上（与未防御基线同一内存后端、检索配置和冻结判断器）测量读取端 reranker 防御的副作用。作者报告该防御使核心准确率下降 4.4 点（95% CI [-9.0, 0.05]，bootstrap；McNemar p=0.064），在无攻击对话中 33.6% 的合法记忆被误隔离，单个对话最多出现 106 次误隔离，令牌开销为 2.7%；作者指出这一准确率变化处于分辨率边缘，机制性代价更显著\cite{W0312}。本文综合认为，防御评估若只报攻击成功率而不报良性成本，将系统性地高估防御可用性；W0312 与 W0178 的对照进一步表明，来源感知的评分机制可能在防投毒与保效用之间取得比纯读取端过滤更好的平衡，但这一推断需在统一基准上验证【待核验】。

\begin{figure}[t]
\centering
\begin{tikzpicture}[x=1.35cm,y=1.05cm,
  phase/.style={rectangle,rounded corners,draw,fill=blue!8,font=\scriptsize,align=center,minimum width=2.2cm},
  risk/.style={font=\tiny,color=red!60!black,align=center}]
\node[phase] (write) at (0,0) {写入\\Write};
\node[phase] (store) at (2,0) {存储\\Store};
\node[phase] (retrieve) at (4,0) {检索\\Retrieve};
\node[phase] (use) at (6,0) {使用\\Use};
\node[phase] (forget) at (2,-1.4) {遗忘/回滚\\Forget};
\node[phase] (share) at (4,-1.4) {共享/传播\\Share};
\draw[->,>=stealth,thick] (write) -- (store);
\draw[->,>=stealth,thick] (store) -- (retrieve);
\draw[->,>=stealth,thick] (retrieve) -- (use);
\draw[->,>=stealth,thick] (store) -- (forget);
\draw[->,>=stealth,thick] (store) -- (share);
\draw[->,>=stealth,dashed,thick] (use) to[bend left=25] (write);
\node[risk] at (0,0.75) {错误写入\\投毒注入};
\node[risk] at (2,0.75) {摘要失真\\来源丢失};
\node[risk] at (4,0.75) {过时信息\\检索错配};
\node[risk] at (6,0.75) {错误使用\\过度泛化};
\node[risk] at (2,-2.05) {误删/残留衍生};
\node[risk] at (4,-2.05) {跨智能体污染};
\node[risk] at (3.2,1.45) {使用反馈驱动再写入（红色为本文关注的关键风险）};
\end{tikzpicture}
\caption{记忆生命周期与关键风险示意（概念图，依据综述文献的生命周期框架
归纳\cite{W0039}；红色标注各阶段风险，详见第~\ref{sec:update}~节）。}
\label{fig:lifecycle}
\end{figure}

% evaluation：内容待 G4/G5 基于证据生成；每条核心陈述将关联 claim_id 注释。
\section{评测与比较}

（本节内容待生成。）

% challenges：内容待 G4/G5 基于证据生成；每条核心陈述将关联 claim_id 注释。
\section{挑战与研究方向}

（本节内容待生成。）

% conclusion：生成于阅读卡+证据；引用检查通过


\section{结论}

现有长期记忆系统的基础机制可归纳为三条技术脉络。第一条是以生成式智能体架构为代表的“记忆流+检索+反思”范式，通过自然语言记录经验，并在检索时综合相关性、新近性与重要性，再以反思模块生成高层推断\cite{W0001}。第二条是虚拟上下文管理路线，以 MemGPT 为代表，将主上下文划分为只读系统指令、可写工作上下文与 FIFO 消息队列，并借助函数调用实现对外部上下文的自主读写与驱逐\cite{W0002}。第三条是读写分离与经验提取路线：RET-LLM 通过微调生成记忆 API 调用，在不修改 LLM 架构的前提下实现显式存储与检索\cite{W0680}；ExpeL 通过自主收集经验并提取跨任务洞察，将回忆对象从原始经验扩展为自然语言知识\cite{W0618}。此外，MemoryBank 引入艾宾浩斯遗忘曲线启发的记忆更新机制，以记忆强度衰减模拟人类记忆的动态性\cite{W0003}；Reflexion 则以语言化自我反思构成情景经验记忆，作者报告其在 HumanEval 编码基准上以自反策略达到 91\% pass@1\cite{W0004}。综合来看，基础系统层已形成“写入—存储—检索—更新”的通用骨架，但在写入时机、记忆粒度和更新策略上仍缺乏统一约定\cite{W0001,W0002,W0680,W0618}。

\subsection{RQ1：长期记忆与相关概念的边界是什么？}

本文的操作定义将长期记忆界定为：由智能体显式写入并跨会话持久化的信息状态。这一边界排除了三类相关但不同的机制：上下文扩展仅提供会话内召回，不构成跨会话的持久写入\cite{W0039}；静态检索增强的外部知识库若无智能体写入行为且不随交互演化，不进入记忆生命周期\cite{W0039}；参数化知识在推理时无法被精确更新，也难以携带个体交互轨迹\cite{W0002,W0618}。据此，综述覆盖的工作以"写入—存储—检索—使用—更新"的显式闭环为共同特征，而机制形态（记忆流、图结构、分层摘要、经验规则）与记忆类型（事实、事件、任务状态、程序性经验）构成设计空间的两条主轴\cite{W0140,W0039}。

\subsection{RQ2：组织与检索机制如何影响记忆系统的性能？}

组织与检索机制构成了长期记忆系统性能分化的关键因素。时间知识图谱类系统（如 Zep）通过情节子图、语义实体子图与社区子图三层结构组织信息，并引入图感知重排器以提升检索质量\cite{W0684}。Mem++ 采用非破坏性组织策略，在写入时完整保存带日期与作者的文档，读取时按问题时间上界筛选，并在冲突场景中保留版本而非覆盖\cite{W0028}。分层图记忆系统 HiGMem 则通过 LLM 引导的证据集裁剪，在 LoCoMo10 基准上将平均检索轮数从 A-Mem 的 99.84 降至 8.09，同时取得更高精确率与相近召回率\cite{W0030}。扩散激活机制（Synapse）通过融合情景与语义记忆，在 LoCoMo 上取得加权平均 F1 40.5，超过 Zep 的 39.7 与 AriGraph 的 33.7\cite{W0099}。然而，组织复杂性的增益并非线性：Mem0 结合图记忆的整体得分仅比基础版高约 2\%\cite{W0682}；一项图结构与扁平向量基线的对比研究显示，图结构记忆的 token F1 为 0.417，低于扁平向量基线的 0.468\cite{W0034}。这表明，在特定评测条件下，组织机制的优势高度依赖于任务与检索设置，不能简单假定结构化组织必然优于扁平存储\cite{W0034,W0682}。

\subsection{RQ3：更新与遗忘机制如何影响长期一致性？}

更新与遗忘机制对长期一致性具有双重作用。在正向更新方面，检索驱动再巩固框架 REALM 将插入、合并与遗忘统一为单一整合阶段，在 LoCoMo 上平均准确率达 75.97\%，比最强基线高 7.17 个百分点\cite{W0012}。SodaMem 通过 supersession 机制在新事实写入时主动失效旧事实，并以软时间窗口而非硬过滤处理时序查询\cite{W0222}。在遗忘治理方面，Memory Worth 估计在 10,000 回合后与真实效用的斯皮尔曼秩相关达 ρ=0.89±0.02，显示在线价值估计具备可行性\cite{W0199}。然而，遗忘也可能削弱系统能力：选择性遗忘一次应用移除 9.8\% 节点后，judged correctness 下降 1.6 点\cite{W0034}；记忆规范化研究显示，全规范化记忆的情感一致性改善在多重比较校正后不再显著，事实漂移在任何校正水平下均不显著，作者明确将其定位为探索性发现\cite{W0077}。这表明更新与遗忘机制在提升一致性的同时，可能引入信息损失风险，需在价值估计与保留策略之间谨慎权衡\cite{W0034,W0077,W0199}。

\subsection{RQ5：现有评测框架能否揭示长期记忆的真实能力？}

现有评测框架已在覆盖维度上取得显著进展，但对长期记忆真实能力的刻画仍不充分。LoCoMo 将对话规模扩展至约 16K tokens、588.2 turns、27.2 sessions，揭示了长上下文 LLM 与 RAG 在时间与因果动态理解上的显著落后\cite{W0005}。LongMemEval 明确将知识更新与弃答纳入五大核心能力\cite{W0006}。Memora 则进一步指出，LoCoMo 中 94\% 的评估问题只需不超过 2 个此前会话的证据，LongMemEval 中该比例为 85\%，显示既有基准偏重浅层检索而缺乏记忆整合压力\cite{W0171}。成本感知评测表明，长期记忆可将依赖任务成功率从无记忆的 0.00 提升至 0.55–1.00，但嵌入检索在事实变化时表现不稳定（0.30–0.95），且代理仅在 55\% 的时间对正确检索值采取行动\cite{W0031}。操作级幻觉评测 HaluMem 则揭示，现有系统倾向于在抽取和更新阶段产生并累积幻觉，再向问答阶段传播\cite{W0261}。这些证据共同指向一个结论：当前评测正向更细粒度的记忆操作和可靠性维度演进，但“记忆何时真正帮助、何时引入危害”仍缺乏系统化的成本—收益度量\cite{W0031,W0171,W0261}。

\subsection{RQ4：可靠性与安全风险如何约束记忆系统的设计？}

可靠性与安全风险已成为长期记忆系统不可回避的设计约束。在非恶意可靠性方面，来源感知系统 Agent Zero Memory 通过引用锁纪律在 LongMemEval 上达 95.60\%，较此前最强系统提升 +0.73 分\cite{W0010}。多智能体错误传播研究显示，幻觉检测率从 Stage 1 的 72.0\% 单调下降至 Stage 4 的 50.9\%，末端检查仅改善 2.3 个百分点\cite{W0688}。信息压缩的统一速率—失真视角进一步指出，查询无关压缩的有效每查询预算会被查询分布熵削减，导致快照式方法在查询分布变化时性能塌陷\cite{W0207}。在安全投毒方面，MemPoison 显示选择性记忆机制可被绕过，攻击成功率最高达 0.95\cite{W0319}；而记忆投毒防御在良性流量中亦有代价——读取端防御使核心准确率下降 4.4 点，33.6\% 的合法记忆被误隔离\cite{W0312}。安全综述进一步将长期记忆界定为“可写、跨会话持久”的新安全基底，并指出保留策略可能意外保留对抗性条目而丢弃合法条目\cite{W0039}。综合而言，可靠性与安全约束要求记忆系统在来源验证、写入控制与良性效用之间寻找平衡，单一环节的强化不足以应对传播链上的综合风险\cite{W0010,W0312,W0039,W0688}。



\bibliographystyle{unsrtnat}
\bibliography{references}

\end{document}
